\section{Data and Lexicon}

\label{sec:data}

\subsection{The UIT-ViHSD corpus}

\label{sec:vihsd}

We use UIT-ViHSD \cite{luu2021vihsd}, Vietnamese social-media comments annotated
with three mutually exclusive labels: \clean{} (no offensive content), \off{}
(offensive language not targeting a protected group as hate speech) and \hate{}
(hateful or seriously derogatory). The \off{}/\hate{} distinction depends on target
and intensity rather than on a particular word, so a lexical rule engine is
insufficient on its own.

After cleaning, the corpus has 33{,}398 comments; every system reads the same fixed
split, whose SHA-256 digest is stored and re-checked per run. Training: 24{,}046
(2{,}556 \hate{}, 1{,}605 \off{}, 19{,}885 \clean{}); validation 2{,}672
(270/212/2{,}190); test 6{,}680 (688/444/5{,}548).

The test split is imbalanced: 83.05\% \clean{}, 10.30\% \hate{} and 6.65\% \off{},
so accuracy is a poor headline metric---a constant \clean{} predictor already reaches
0.8305---and macro-F1 is primary, always with per-class F1. The 6{,}680 test rows
contain 6{,}571 distinct texts (109 duplicated rows), kept because removing them would
change the frozen test set; one appears twice among the 15 regressions, which thus
cover 14 distinct texts.

\subsection{The hate-speech lexicon}

\label{sec:lexicon}

The rule engine is driven by a Vietnamese hate-speech lexicon of \textbf{668 entries},
mined and labelled from the \emph{training} portions of UIT-ViHSD; R-UIT-ViHSD
\cite{bontoriTdtu}; and ViCM \cite{nguyen2026vicm}. The three training splits yield
1{,}062 candidate terms (first attested 257/665/140 in UIT-ViHSD/ViCM/R-UIT-ViHSD); manual labelling and
filtering reduce these candidates to the 668 entries below. No evaluation split is sampled; each entry records
its supporting corpora (414 of 668 occur in all three).

An exact-text cross-corpus audit found 895 of 6{,}680 UIT-ViHSD test rows (13.4\%)
with a matching text in auxiliary corpus material used for lexicon development---620
distinct test texts also in the Bontori/R-UIT-ViHSD training partition---a potential
source of lexicon-construction leakage. To assess its effect, we rebuilt the lexicon
excluding those texts (671 entries vs. 668 canonical). Main results use the canonical
lexicon; leak-free results are in Sect.~\ref{sec:leak-free-sensitivity}.

Each entry stores the surface term and canonical form, a label (\hate{} or \off{}) with
per-label severity, a training-estimated confidence, a context-required flag, and a
review status (\emph{verified} or \emph{auto}), retained in the evidence record.

The lexicon is deliberately not a list of slurs: 551 of 668 (82.5\%) are multi-word
constructions (\emph{mẹ cha}, \emph{con mẹ}, \emph{ngu hơn}) and 524 (78.4\%) are
context-required, since Vietnamese insults work by comparison and implication rather
than a single taboo word. Only 74 entries are \off{} against 594 \hate{}, so the engine
is weak on \off{} and fusion keeps that channel mostly neural. The lexicon is
\emph{editable}, \emph{deterministic} and \emph{inspectable} (each score decomposes
into term, match type and context factor).

\paragraph{Ethics of examples.}
All quoted comments are public UIT-ViHSD test comments. Because they contain profanity
and slurs, we mask offending characters with \texttt{*} in every table; evaluation
runs on unmasked text. No user identifiers are reproduced, and names inside comments
are part of the public corpus text.